\section{Related Work}
\label{sec:related}

\paragraph{RNA structure prediction and inverse folding.}
Recent models have advanced RNA structure prediction across tertiary and secondary levels. For tertiary structure, traditional methods like DRFold \citep{li2023integrating}  and DeepFoldRNA \citep{pearce2022novo} combine primary/secondary structure information with energy functions for prediction and optimization. AlphaFold2-inspired approaches including NuFold \citep{kagaya2025nufold}, trRoseTTARNA \citep{wang2023trrosettarna} and RhoFold+ \citep{shen2024accurate} employ EvoFormer architectures to predict structures from sequence and secondary structure data \citep{jumper2021alphafold2}. RoseTTAFoldNA \citep{baek2024accurate} and AlphaFold 3 \citep{abramson2024accurate} extend prediction capabilities to protein-nucleic acid complexes and multi-component assemblies. For secondary structure, traditional methods RNAfold/ViennaRNA \citep{lorenz2011viennarna} calculate partition functions and base-pairing probabilities through minimum free energy optimization. SPOT-RNA \citep{singh2019rna} uses 2D CNNs with transfer learning for non-canonical base pairs, E2EFold \citep{chen2020rna} predicts contact matrices, and RFold \citep{tan2024deciphering} employs bidirectional decomposition to constrain cross-pairing interactions. Current inverse folding methods predominantly use GNN-based structural encoders, including RNAinformer \citep{patil2024towards}, R3Design \citep{tan2025r3design}, RDesign \citep{tanrdesign}, RhoDesign \citep{wong2024deep}, RiFold \citep{liu2025sentences}, AlignIF \citep{wang2025alignment}, and gRNAde \citep{joshi2025grnade}. RiboDiffusion \citep{huang2024ribodiffusion} alternatively employs structure-conditioned discrete diffusion models. Notable advances include RhoDesign's \citep{wong2024deep} experimental validation, RiFold's language modeling framework \citep{liu2025sentences}, and gRNAde's \citep{joshi2025grnade} extension from single-state to multi-state design.

\paragraph{Reinforcement Learning and Preference Optimization for Inverse Folding.}
RL and preference optimization are increasingly used to steer inverse folding toward structural fidelity and stability. Structure feedback has already been used for protein inverse folding optimization \citep{xu2025proteininversefoldingstructure}. ProteinZero explores online RL for self-improvement with proxy rewards \citep{wang2025proteinzero}, while ResiDPO introduces residue-level Direct Preference Optimization (DPO) using structure-based rewards to improve designability \citep{xue2025improving}. Diversity-regularized DPO has been applied to peptide inverse folding \citep{park2024improving}. PLM-RL combined RL with the Protein Language Models \citep{cao2025supervisionexplorationdoesprotein}. RL has also been combined with diffusion models: RL-DIF optimizes structure-conditioned categorical diffusion \citep{ektefaie2024reinforcement}, and DRAKES backpropagates rewards through discrete diffusion trajectories for DNA/protein design \citep{wang2024fine}. Search-based strategies like ProtInvTree use reward-guided tree search \citep{liu2025protinvtree}, and lightweight, gradient-free finetuning for discrete sequences is explored in GLID$^2$E \citep{cao2025glid}. Finally, EnerBridge-DPO integrates energy-based preferences via Markov bridges with DPO to bias designs toward low-energy sequences \citep{rong2025enerbridge}. To the best of our knowledge, none of existing methods has demonstrated the effectiveness of RLPF on RNA.